\section{Subgroups}

\subsection{Definitions and Examples}

\begin{exercise}[2.1.1]
    In each of (a) -- (e) prove that the specified subset is a subgroup of the given group.
    \begin{subexercises}
        \item The set of complex numbers of the form $a + ai$, $a \in \r$ (under addition).
        \item The set of complex numbers of absolute value $1$, i.e., the unit circle in the complex plane (under multiplication).
        \item For fixed $n \in \z^+$, the set of rational numbers whose denominators divide $n$ (under addition).
        \item For fixed $n \in \z^+$, the set of rational numbers whose denominators are relatively prime to $n$ (under addition).
        \item The set of nonzero real numbers whose square is a rational number (under multiplication).
    \end{subexercises}
\end{exercise}

\begin{sol}
    Let $G$ be the group and $H$ be the subset in question.
    \begin{subexercises}
        \item Since $0 + 0i \in H$, then $H$ is nonempty. For $a + ai, b + bi \in H$,
        \[(a + ai) - (b + bi) = (a - b) + (a - b)i \in H\]
        so that $H \leq G$.
        \item Since $\abs 1 = 1$, then $1 \in H$ so that it is nonempty. For $z, w \in H$, we first note that $|w\inv| = 1/|w| = 1/1 = 1$. Then
        \[|zw\inv| = |z||w\inv| = 1 \cdot 1 = 1\]
        so that $zw\inv \in H$, hence $H \leq G$.
        \item Fix $n \in \zp$. Since $0 = 0/k$ for some $k$ such that $k \mid n$, then $0 \in H$ so that $H$ is nonempty. Let $a/b, c/d \in H$ such that $b \mid n$ and $d \mid n$. Then there exist $x, y \in \zp$ such that $bx = n$ and $dy = n$. Then
        \[\frac ab - \frac cd = \frac{ax - cy}{n}.\]
        Reducing this fraction still yields a denominator that divides $n$. Then $H \leq G$.
        \item Fix $n \in \zp$. Since $0 = 0/k$ for some $k$ such that $(k, n) = 1$, then $0 \in H$ so that $H$ is nonempty. Let $a/b, c/d \in H$ such that $(a, c) = 1$, $(b, n) = 1$, and $(d, n) = 1$. Then
        \[\frac ab - \frac cd = \frac{ad - bc}{bd}\]
        If $(bd, n) \neq 1$, then there exists some prime $p$ such that $p \mid bd$ and $p \mid n$. Then $p \mid b$ or $p \mid d$, contradicting our assumption. Hence, $(bd, n) = 1$, so that $H \leq G$.
        \item Since $1 = 1^2 = 1/1$, then $1 \in H$ so that $H$ is nonempty. For $x, y \in H$, then $x^2, y^2 \in \q$. Then
        \[\left(\frac xy\right)^2 = \frac{x^2}{y^2} \in \q\]
        so that $x/y \in H$, hence $H \leq G$. \qh
    \end{subexercises}
\end{sol}

\begin{exercise}[2.1.2]
    In each of (a) -- (e) prove that the specified subset is \emph{not} a subgroup of the given group.
    \begin{subexercises}
        \item The set of $2$-cycles in $S_n$ for $n \ge 3$.
        \item The set of reflections in $D_{2n}$ for $n \ge 3$.
        \item For $n$ a composite integer $>1$ and $G$ a group containing an element of order $n$, the set $\set{x \in G \mid |x| = n} \cup \set{1}$.
        \item The set of (positive and negative) odd integers in $\z$ together with $0$.
        \item The set of real numbers whose square is a rational number (under addition).
    \end{subexercises}
\end{exercise}

\begin{sol}
    Let $G$ be the group and $H$ be the subset in question.
    \begin{subexercises}
        \item For $n = 3$, then $(1\ 2), (1\ 3) \in H$, but $(1\ 2)(1\ 3) = (1\ 3\ 2) \notin H$.
        \item For $n = 3$, then $s, sr \in H$, but $s(sr) = r \notin H$.
        \item Let $n = ab$ for some $a, b \in \zp$ with $1 < a, b < n$. Let $x \in G$ with $\abs x = n$. Then $x \in H$ so that $x^a \in H$, but $(x^a)^b = x^n = 1$ implies that $|x^a| \leq b$.
        \item Observe that $1 \in H$, but $1 + 1 = 2 \notin H$.
        \item Observe that $\sqrt 2, \sqrt 3 \in H$, but $(\sqrt 2 + \sqrt 3)^2 \notin \q$. \qh
    \end{subexercises}
\end{sol}

\begin{exercise}[2.1.3]
    Show that the following subsets of the dihedral group $D_8$ are actually subgroups:
    \begin{subexercises}
        \item $\set{1, r^2, s, sr^2}$
        \item $\set{1, r^2, sr, sr^3}$
    \end{subexercises}
\end{exercise}

\begin{solenum}
    \item 
    \[
    \begin{array}{c|cccc}
        \circ & 1 & r^2 & s & sr^2 \\
        \hline
        1 & 1 & r^2 & s & sr^2 \\
        r^2 & r^2 & 1 & sr^2 & s \\
        s & s & sr^2 & 1 & r^2 \\
        sr^2 & sr^2 & s & r^2 & 1
    \end{array}
    \]
    \item 
    \[
    \begin{array}{c|cccc}
        \circ & 1 & r^2 & sr & sr^3 \\
        \hline
        1 & 1 & r^2 & sr & sr^3 \\
        r^2 & r^2 & 1 & sr^3 & sr \\
        sr & sr & sr^3 & 1 & r^2 \\
        sr^3 & sr^3 & sr & r^2 & 1
    \end{array}
    \]
    Both tables show closure, since no product is an element outside the subset.
\end{solenum}

\begin{exercise}[2.1.4]
    Give an explicit example of a group $G$ and an infinite subset $H$ of $G$ that is closed under the group operation but is not a subgroup of $G$.
\end{exercise}

\begin{sol}
    Let $G = \z$ and $H = \zp$. Then for any $m, n \in \zp$, we have $m + n \in \zp$ but $m - n \notin \zp$ when $m < n$. Hence, $H$ is not a subgroup of $G$.
\end{sol}

\begin{exercise}[2.1.5]
    Prove that $G$ cannot have a subgroup $H$ with $|H| = n - 1$, where $n = |G| > 2$.
\end{exercise}

\begin{sol}
   If $H \leq G$ with $|H| = n - 1$, then there exists $k \in \zp$ such that $k(n - 1) = n$ by Lagrange's Theorem. This implies $kn - k = n$, or equivalently $n = k/(k - 1)$. Since $k = 2$ implies $n = 2$ and $k \geq 3$ implies $n \notin \zp$, then there is no such subgroup $H$.
\end{sol}

\begin{spex}[2.1.6]
    \label{def:torsion subgroup}
    Let $G$ be an Abelian group. Prove that $\set{ g \in G \mid |g| < \infty }$ is a subgroup of $G$ (called the \textbf{torsion subgroup} of $G$). Give an explicit example where this set is not a subgroup when $G$ is non-Abelian.
\end{spex}

\begin{sol}
    Let the torsion subgroup be denoted as $\tor(G)$. Since $|1| = 1$, then $1 \in \tor(G)$. If $g, h \in \tor(G)$ such that $|g| = m$ and $|h| = n$, we have
    \[(gh\inv)^{mn} = g^{mn}(h\inv)^{mn} = (g^m)^n((h^n)\inv)^m = 1^n(1\inv)^m = 1,\]
    where the first equality follows from $G$ being Abelian. Then $|gh\inv| \leq mn$, hence $gh\inv \in \tor(G)$, and $\tor(G) \leq G$.

    Consider $G = \gl_2(\r)$ where
    \[A = 
    \begin{pmatrix}
        -1 & 0 \\
        0 & 1
    \end{pmatrix}, \quad
    B = 
    \begin{pmatrix}
        -1 & 1 \\
        0 & 1
    \end{pmatrix}.\]
    Note that $|A| = |B| = 2$ so that $A, B \in \tor(G)$. However,
    \[AB = 
    \begin{pmatrix}
        1 & -1 \\
        0 & 1
    \end{pmatrix}.\]
    By induction, we may show that
    \[(AB)^n = 
    \begin{pmatrix}
        1 & -n \\
        0 & 1
    \end{pmatrix}\]
    for all $n \in \zp$. Hence, $|AB| = \infty$, so $AB \notin \tor(G)$.
\end{sol}

\begin{exercise} [2.1.7]
    Fix some $n \in \z$ with $n>1$. Find the \defref{torsion subgroup} of $\z \times (\z/n\z)$. Show that the set of elements of infinite order together with the identity is \emph{not} a subgroup of this direct product.
\end{exercise}

\begin{sol}
    The only finite order element in $\z$ is 0, while $\intmod$ is finite. Then
    \[\tor(\z \times \intmod) = \set{(0, \bar x) \mid \bar x \in \intmod}.\]
    Consider $H = \set{(x, \bar y) \in \z \times \intmod \mid x \neq 0} \cup \set{(0, \bar 0)}$. Then $H$ contains the identity and all elements of infinite order. However, $(1, \bar 1), (-1, \bar 0) \in H$, but $(0, \bar 1) = (1, \bar 1) + (-1, \bar 0) \notin H$, since it has finite order. Hence, $H$ is not a subgroup of $\z \times \intmod$.
\end{sol}

\begin{exercise} [2.1.8]
    Let $H$ and $K$ be subgroups of $G$. Prove that $H \cup K$ is a subgroup if and only if either $H \subseteq K$ or $K \subseteq H$.
\end{exercise}

\begin{solitem}
    \item[\rightimp] Suppose $H \cup K$ is a subgroup. If $H \subseteq K$, then we are done. Suppose that $H \nsubseteq K$. Then there exists some $h \in H$ such that $h \notin K$. Since $H \cup K$ is a subgroup, then for any $k \in K$, we have $hk \in H \cup K$.
    
    Observe that $k\inv \in K$ since $K \leq G$. If $hk \in K$, then $h = (hk)k\inv \in K$, contradicting our assumption. It must be that $hk \in H$. Since $h \in H$, then $h\inv \in H$ as $H \leq G$. Then $h\inv(hk) = k \in H$, hence $K \subseteq H$.
    \item[\leftimp] If $H \subseteq K$, then $H \cup K = K \leq G$. Similarly, if $K \subseteq H$, then $H \cup K = H \leq G$.
\end{solitem}

\begin{spex}[2.1.9]
    Let $G = \gl_n(F)$, where $F$ is any field. Define
    \[\speclin_n(F) = \set{ A \in \gl_n(F) \mid \det(A) = 1 }\]
    (called the \textbf{special linear group}). Prove that $\speclin_n(F) \le \gl_n(F)$.
\end{spex}

\begin{sol}
    Since $\det(I_n) = 1$, then $I_n \in \speclin_n(F)$. Suppose $A, B \in \speclin_n(F)$. Then
    \[\det(AB\inv) = \det(A)\det(B\inv) = \frac{\det(A)}{\det(B)} = \frac 11 = 1\]
    so that $AB\inv \in \speclin_n(F)$. It follows that $\speclin_n(F) \leq \gl_n(F)$.
\end{sol}

\begin{exercise} [2.1.10]
    \begin{subexercises}
        \item Prove that if $H$ and $K$ are subgroups of $G$ then so is their intersection $H \cap K$.
        \item Prove that the intersection of an arbitrary nonempty collection of subgroups of $G$ is again a subgroup of $G$ (do not assume the collection is countable).
    \end{subexercises}
\end{exercise}

\begin{sol}
    We prove (b) and note that (a) is a special case with $I = \set{1, 2}$ and $H_1 = H$, $H_2 = K$. Let $I$ be an indexing set, and consider the collection of subgroups $\set{ H_i }_{i \in I}$ of $G$. Let
    \[H = \bigcap_{i \in I} H_i\]
    Since $1 \in H_i$ for all $i \in I$, then $1 \in H$. Suppose $a, b \in H$. Then $a, b \in H_i$ for all $i \in I$. Since each $H_i$ is a subgroup of $G$, then $ab\inv \in H_i$ for all $i \in I$. Hence, $ab\inv \in H$, and $H \leq G$.
\end{sol}

\begin{exercise} [2.1.11]
    Let $A$ and $B$ be groups. Prove that the following sets are subgroups of the direct product $A \times B$:
    \begin{subexercises}
        \item $\set{ (a, 1) \mid a \in A }$
        \item $\set{ (1, b) \mid b \in B }$
        \item $\set{ (a, a) \mid a \in A }$, where we assume $B = A$ (called the diagonal subgroup)
    \end{subexercises}
\end{exercise}

\begin{sol}
    Let $C$ be the set in each part.
    \begin{subexercises}
        \item Since $1 \in A$, then $(1, 1) \in C$ so that $C$ is nonempty. Suppose $(a_1, 1), (a_2, 1) \in C$. Then
        \[(a_1, 1)(a_2, 1)\inv = (a_1, 1)(a_2\inv, 1) = (a_1a_2\inv, 1) \in C\]
        since $a_1a_2\inv \in A$. Hence, $C \leq A \times B$.
        \item Since $1 \in B$, then $(1, 1) \in C$ so that $C$ is nonempty. Suppose $(1, b_1), (1, b_2) \in C$. Then
        \[(1, b_1)(1, b_2)\inv = (1, b_1)(1, b_2\inv) = (1, b_1b_2\inv) \in C\]
        since $b_1b_2\inv \in B$. Hence, $C \leq A \times B$.
        \item Since $1 \in A$, then $(1, 1) \in C$ so that $C$ is nonempty. Suppose $(a_1, a_1), (a_2, a_2) \in C$. Then
        \[(a_1, a_1)(a_2, a_2)\inv = (a_1, a_1)(a_2\inv, a_2\inv) = (a_1a_2\inv, a_1a_2\inv) \in C\]
        since $a_1a_2\inv \in A$. Hence, $C \leq A \times B$. \qh
    \end{subexercises}
\end{sol}

\begin{exercise} [2.1.12]
    Let $A$ be an Abelian group and fix some $n \in \z$. Prove that the following sets are subgroups of $A$.
    \begin{subexercises}
        \item $\set{ a^n \mid a \in A }$.
        \item $\set{ a \in A \mid a^n = 1 }$.
    \end{subexercises}
\end{exercise}

\begin{sol}
    Let $B$ be the sets in question.
    \begin{subexercises}
        \item Since $1 = 1^n$, then $1 \in B$ so that $B$ is nonempty. Suppose $a^n, b^n \in B$. Then
        \item \[(a^n)(b^n)\inv = a^n b^{-n} = (ab\inv)^n \in B\]
        where the last equality follows from the fact that $A$ is Abelian. Hence, $B \leq A$.
        \item Since $1^n = 1$, then $1 \in B$ so that $B$ is nonempty. Suppose $a, b \in B$. Then
        \[(ab\inv)^n = a^n (b\inv)^n = a^n (b^n)\inv = 1 \cdot 1\inv = 1\]
        where the second equality follows from the fact that $A$ is Abelian. Hence, $B \leq A$. \qh
    \end{subexercises}
\end{sol}

\begin{exercise} [2.1.13]
    Let $H$ be a subgroup of the additive group of rational numbers with the property that $1/x \in H$ for every nonzero element $x$ of $H$. Prove that $H = 0$ or $\q$.
\end{exercise}

\begin{sol}
    Let $H \leq \q$. Then $0 \in H$. If no other elements are in $H$, then $H = 0$, and we are done. Suppose there exists some nonzero $a/b \in H$ where $a, b \in \z$ and $b \neq 0$. Moreover, we may assume without loss of generality that $a/b > 0$, for otherwise we may take $-a/b$ since $H$ contains additive inverses. We now use $a/b$ to generate $\q$.

    Since $a/b \in H$, then $b(a/b) = a \in H$ so that $1/a \in H$. Then $a(1/a) = 1 \in H$. From this, we can see that $\z \subseteq H$. Now let $c/d \in \q$ for some $c, d \in \z$ with $d \neq 0$. Since $1 \in H$, then $d(1) = d \in H$ so that $1/d \in H$. Then $c(1/d) = c/d \in H$. Since $c/d$ was arbitrary, then $\q \subseteq H$. Hence, $H = \q$.
\end{sol}

\begin{exercise} [2.1.14]
    Show that $\set{ x \in D_{2n} \mid x^2 = 1 }$ is \emph{not} a subgroup of $D_{2n}$ (here $n \ge 3$).
\end{exercise}

\begin{sol}
    Let $H$ be the set in question. All reflections are in $H$. In particular, $s$ and $sr$ are in $H$, but $r = s(sr) \notin H$ since $\abs r = n > 2$.
\end{sol}

\begin{exercise} [2.1.15]
    Let $H_1 \subseteq H_2 \subseteq \cdots$ be an ascending chain of subgroups of $G$. Prove that
    \[\bigcup_{i = 1}^\infty H_i\]
    is a subgroup of $G$.
\end{exercise}

\begin{sol}
    Let $H$ be the union in question. From $H_i \leq G$ for each $i \in \zp$, then $1 \in H_i$ so that $1 \in H$. Let $a, b \in H$ so that there exists $m, n \in \zp$ where $a \in H_m$ and $b \in H_n$. Put $k = \max(m, n)$. It follows that $H_m \subseteq H_k$ and $H_n \subseteq H_k$, so that $a, b \in H_k$. Since $H_k \leq G$, then $ab\inv \in H_k$, hence $ab\inv \in H$. Therefore, $H \leq G$.
\end{sol}


\begin{exercise} [2.1.16]
    Let $n \in \z^+$ and let $F$ be a field. Prove that the set $\set{ (a_{ij}) \in \gl_n(F) \mid a_{ij} = 0 \text{ for all } i > j }$ is a subgroup of $\gl_n(F)$ (called the group of upper triangular matrices).
\end{exercise}

\begin{sol}
    Let $\ut_n(F)$ denote the set of $n \times n$ upper triangular matrices with entries from $F$. Observe that $I_n$ contains $0$'s everywhere except the diagonal, hence $I_n \in \ut_n(F)$ so that $\ut_n(F)$ is nonempty. Suppose $A, B \in \ut_n(F)$, and consider $C = AB$. Then the $(i, j)$-th entry of $C$ is given by
    \[c_{ij} = \sum_{k = 1}^n a_{ik}b_{kj}\]
    Consider the case when $i > j$, which is below the main diagonal.
    \begin{itemize}
        \item If $i > k$, then $a_{ik} = 0$ since $A$ is upper triangular.
        \item If $k > j$, then $b_{kj} = 0$ since $B$ is upper triangular.
    \end{itemize}
    Hence, for all $k$, either $a_{ik} = 0$ or $b_{kj} = 0$, so that $c_{ij} = 0$. It follows that $C \in \ut_n(F)$.

    Let $A \in \ut_n(F)$. Since $A \in \gl_n(F)$, we know that the entries on the main diagonal are all non-zero, for otherwise $A$ is not invertible. By back-substitution, we may solve the matrix equation $AX = I_n$ for $X$, and the solution will also be upper triangular. Hence, $A\inv \in \ut_n(F)$.
\end{sol}

\begin{exercise} [2.1.17]
    Let $n \in \z^+$ and let $F$ be a field. Prove that the set $\set{ (a_{ij}) \in \gl_n(F) \mid a_{ij} = 0 \text{ for all } i > j,\ \text{and } a_{ii} = 1 \text{ for all } i }$ is a subgroup of $\gl_n(F)$.
\end{exercise}

\begin{sol}
    Let $H$ be the set in question. Since the diagonal of $I_n$ is only 1's, then $I_n \in H$. Suppose $A, B \in H$. By the previous exercise, $A, B \in \ut_n(F)$, so it remains to show that $a_{ii} = b_{ii} = 1$ for all $1 \leq i \leq n$. Then
    \[(AB)_{ii} = \sum_{k = 1}^n a_{ik}b_{ki}\]
    Since $a_{ik} = 0$ for $k < i$ and $b_{ki} = 0$ for $k > i$, then the sum degrades to $a_{ii}b_{ii} = 1$. Then $H$ is closed under multiplication. Moreover, suppose $D \in \ut_n(F)$ such that $DA = I_n$. Then
    \[1 = (DA)_{ii} = d_{ii}a_{ii}\]
    where we use the above to collapse the $ii$-th term of $DA$. Then $d_{ii} = 1$, hence $D \in H$, and $H$ is closed under inverses. Hence, $H \leq \gl_n(F)$.
\end{sol}